% PDF-derived LaTeX transcription. Original manuscript source was unavailable.
% Published PDF SHA-256: a350ed15bef5a283c34f37716ba1c3a9582ec1cffb622f4e7d4a65bfb21c1dd5
% Extracted text SHA-256: 385a959e391612ae3a1e7f66d4f84497d86b9399af02aeff22461f0930eb50d0
\documentclass[11pt]{article}
\usepackage[margin=1in]{geometry}
\usepackage{amsmath,amssymb}
\setlength{\emergencystretch}{3em}
\title{What Passing Solver Checks Establish: Periodic Decay, Manufactured Pressure, and Failed Boundary Response in Fluid Verification}
\author{}
\date{}
\begin{document}
\maketitle
\noindent\textbf{PDF-derived LaTeX transcription.} This editable text follows the PDF extraction. Refer to the published PDF for figures, table layouts, and mathematical notation.

% Extracted chunk 1

What Passing Solver Checks Establish:
Periodic Decay, Manufactured Pressure,
and Failed Boundary Response in Fluid Verification

\section*{Abstract}

A numerical acceptance result applies to an observable, a configuration and a stated toler-
ance. We examine the consequences of keeping those conditions explicit across four retained
fluid-verification campaigns. Ten periodic Taylor–Green cases passed their recorded limits;
stationary decay-energy errors gave observed resolution orders of 1.987–2.054, and a matched
moving-frame comparison differed by 0.0868\% in its energy-retention ratio. A bounded residual-
refresh candidate was admitted on four manufactured pressure graphs, although the original
comparison panel retained failed baseline outcomes. In annular Couette startup, independently
reconstructed observables from 721 saved planar states agreed with the original reductions, while
the physical response still failed: maximum surface-torque errors were 5.53–7.06\% of the steady
torque scale, aggregate profile error was 7.327\%, and final profile error was 5.197\%, against
5\% limits. Separate fixed-flow S-bend records showed nonmonotone pressure refinement and a
corrector-induced change comparable to the medium-to-fine change despite very small surface-
flux imbalance. These are configuration-specific tests, not a controlled isolation of one failure
mechanism. Their combined contribution is an empirical account of why periodic accuracy,
algebraic admission, observable reconstruction and physical boundary response require separate
decisions.

\section*{1 Question and contribution}

Fluid software is often examined through a sequence of inexpensive numerical tests before more
demanding application calculations. A passing test can support a useful, narrow conclusion without
qualifying every later use of the implementation. The issue is especially clear when a method passes
periodic decay and algebraic pressure checks, yet misses a physically meaningful velocity or stress
tolerance at an immersed boundary.
This study asks what four retained campaigns establish when their configurations, observables and
evidence levels are kept separate. The first examines periodic vortex decay in a D3Q27 cumulant
implementation. The second assesses a variable-coefficient pressure solve on manufactured con-
nected graphs. The third compares annular Couette startup with a continuum reference and tests
control-surface observables using saved moments. The fourth examines mesh and outer-corrector
sensitivity in a separate OpenFOAM S-bend calculation. Together they retain both positive results
and substantive failures, rather than treating a successful numerical check as transferable approval.
The cumulant lattice Boltzmann formulation and its diagonal corrections are inherited from the
method family described by Geier et al. [ 1]. Our contribution is the specific empirical comparison
of these acceptance layers and the retained-data reconstruction of the Couette observables. We
introduce neither a new collision model nor a general verification theorem. The campaigns differ
1

in geometry, numerical controls and sometimes solver family; they are not matched experiments
that identify one cause of the Couette discrepancy. In particular, the periodic cases enable a local
diagonal correction that the captured Couette configuration disables.

\section*{2 Study design and evidence levels}

The analysis uses original reports, configuration and method records, saved Couette moments, and
surface summaries from the curved-flow calculations. Existing outputs are examined retrospec-
tively. We do not rerun collision and streaming, solve new pressure graphs, or advance new fluid
trajectories. Independent reductions refer to new implementations of calculations on the retained
data; they do not imply independent generation of those data.
Table 1 states the observation contract for each campaign. Case identity and complete source
bindings are retained in a private evidence index. A separate unresolved-check report identifies
missing inputs and the checks they prevent. The scientific results below do not depend on treating
unavailable observations as positive evidence.
Campaign Numerical configuration and endpoint Present evidence and decision scope
Periodic decay CPU; float32 D3Q27 product equi-
librium; diagonal correction enabled;
nine stationary and one moving case
Original reports and inert method defini-
tions; analytical decay and gate arithmetic
checked. No population evolution replay.
Manufactured
pressure
CPU JAX, 64-bit mode disabled;
four connected graphs; fixed residual-
refresh controls
Reported residuals, gradients and admis-
sion limits checked against the source-
defined panel. No new matrix–vector resid-
ual measurement.
Couette startup CPU; float32 evolution; diagonal cor-
rection disabled; 963 lattice, 720 steps
to 28.8 s
Independent moment-based observer reduc-
tion at all 721 saved times; seven explicit
central-plane profiles and 721 norm sum-
maries. Physical profile/stress limits re-
main failed.
S-bend fixed flow OpenFOAM pimpleFoam; three
meshes; 512 steps to 20.48 s; outer-
corrector variants
Original surface summaries and scalar dif-
ferences checked. Mesh, temporal and full-
field qualification remains incomplete.
Table 1: Configurations and evidence levels. A check on retained arithmetic is distinguished from
a new numerical or physical experiment. The three canonical campaigns report CPU execution; no
accelerator speedup is measured.
The samples have different meanings. The Taylor–Green matrix contains ten deterministic cases;
the pressure panel contains four manufactured graphs; the Couette record contains correlated times
from one startup; and the S-bend records contain mesh and iteration variants. We do not pool them
into one statistical sample or infer an empirical probability that a passing residual check predicts
physical accuracy. All tolerances reported here are those specified in the captured test definitions.
Source statements that a panel was fixed before execution are not a substitute for an independently
dated preregistration.
2


\section*{3 Periodic decay: a positive result with a defined frame compar-}

ison

\section*{3.1 Definitions and configuration}

The periodic test uses an x– y Taylor–Green velocity mode with nominal lattice amplitude 0.01, a
product equilibrium, and an enabled local diagonal correction. The nine stationary cases combine
N = 16 , 24, 32 with relaxation times \ensuremath{\tau} = 0 .6, 0.8, 1.0. A tenth case adds a uniform x velocity
U = 0.1 at N = 32, \ensuremath{\tau} = 0.8. The lattice viscosity and wave number are
\ensuremath{\nu}\ensuremath{\ell} = \ensuremath{\tau} \ensuremath{-} 1/2
3 , k = 2\ensuremath{\pi}
N . (1)
The analytical mode amplitude decays as A(n) = A0 exp(\ensuremath{-}2\ensuremath{\nu}\ensuremath{\ell}k2n). Here A0 is the projection of
the actual stored initial velocity onto the test mode, not an assertion that the initialized float32
state has amplitude exactly equal to the nominal input.
The energy observable removes the measured spatial mean of each state:
Ef (n) = 1
2
D
|⃗ u(n) \ensuremath{-} \ensuremath{\langle}⃗ u(n)\ensuremath{\rangle}|2
E
, R (n) = Ef (n)
Ef (0) , R ex(n) = exp( \ensuremath{-}4\ensuremath{\nu}\ensuremath{\ell}k2n). (2)
This is a per-unit-mass lattice fluctuation energy, not the total kinetic energy of the uniform
background. The moving-mode projection is translated by U n in x. The intended amplitude
exponent is 0.5, but integer step counts give slightly different realized decay clocks. We evaluate
the analytical ratios at those actual counts, rather than replacing them with one common terminal
time.
Define the relative energy error as eE = |R/Rex \ensuremath{-}1|. The reported resolution indicator and matched
moving-frame discrepancy are
p16,32 = log[eE(16)/eE(32)]
log 2 , d U =

% Extracted chunk 5


RU =0.1
RU =0
\ensuremath{-} 1

% Extracted chunk 9

. (3)
These definitions make the denominator of the frame comparison explicit. A small dU is a bounded
comparison of two energy-retention ratios; it is not exact Galilean invariance of the implementation.

\section*{3.2 Retained results}

All ten cases satisfy the captured scalar gates, including the 5\% amplitude and energy limits. The
other limits are relative mass error 10\ensuremath{-}6, maximum density departure 0.005, transverse velocity
10\ensuremath{-}5, peak lattice Mach number 0.3, and mode purity at least 0.995. The recorded minimum ac-
cepted 16-to-32 order is 1.0. The observed energy-error orders are 1.986908, 2.000964 and 2.054173
for the three relaxation times. This supports a near-second-order trend over the retained resolution
pairs at their rounded decay clocks, rather than a general asymptotic-order proof.
For the matched N = 32 , \ensuremath{\tau} = 0 .8, n = 65 pair, the stationary energy ratio is 0.3638449309 and
the moving ratio is 0.3641606455. Equation ( 3) gives dU = 0 .0008677175, or 0.08677175\%, below
the captured 2\% limit. The absolute ratio difference is a different quantity and is not the reported
percentage.
The source and report definitions support these analytical and scalar comparisons. They do not
recover the absent initial population arrays, independently reconstruct A0, or reproduce the evo-
lution. Geier et al. [ 1] provide the inherited cumulant and correction framework; their published
3

\ensuremath{\tau} N Steps n e E (\%) Amplitude error (\%)
0.6 16 49 4.70982 2.38332
0.6 24 109 2.10983 1.06054
0.6 32 195 1.18819 0.595871
0.8 16 16 3.44572 1.73797
0.8 24 36 1.53010 0.768005
0.8 32 65 0.860855 0.431366
1.0 16 10 1.34198 0.673330
1.0 24 22 0.577967 0.289403
1.0 32 39 0.323131 0.161726
0.8 32, moving 65 0.774830 0.388210
Table 2: Periodic decay results retained in the original reports. All stationary cases have U = 0 ;
the last case has U = 0 .1. Errors use each case’s actual integer horizon and stored-initial-state
amplitude convention.
Taylor–Green setup has a different aspect ratio and initialization. Our finite test matrix is not a
matched repetition of that published experiment. Their discussion of incomplete local corrections
also precludes promoting this one moving-frame result into a claim of general lattice invariance.

\section*{4 Manufactured pressure: numerical admission is a separate end-}

point
The pressure panel examines connected variable-coefficient graphs with two excluded cells in each
reported shape. Three graphs use 8 \ensuremath{\times} 9 \ensuremath{\times} 10 arrays and coefficient-contrast multipliers 102, 103 and
104; a fourth uses 10\ensuremath{\times}11\ensuremath{\times}12 and multiplier 103. They contain 718 and 1318 fluid cells, respectively.
The source-defined spacing is 0.125 and the four seeds are 20261021– 20261024 in table order. The
contrast control scales coefficients in one half of the third coordinate. Manufactured fields define
the right-hand side, with a pressure gauge on the connected fluid graph.
The retained comparison distinguishes the original float32 right-hand-side reference from a float64
mean-zero-compatible reference. Removing the remaining compatibility defect in float64 changes
the comparison problem slightly, so gradients against both references are reported. An algebraically
well-resolved compatible direct solve does not itself measure a physical pressure field.
For a discrete system Ap = b, the fresh residual and a reference-gradient discrepancy have the forms
rrel = \ensuremath{\parallel}Ap \ensuremath{-} b\ensuremath{\parallel}2,F
\ensuremath{\parallel}b\ensuremath{\parallel}2,F
, g rel = \ensuremath{\parallel}Gp \ensuremath{-} Gpref \ensuremath{\parallel}2
\ensuremath{\parallel}Gpref \ensuremath{\parallel}2
, (4)
where F is the fluid-cell set and G is the source-defined edge-gradient comparison. The sign of b
follows the corresponding operator convention. The report distinguishes freshly applied matrix-free
residuals from recursive or preconditioned stopping histories and from a separate assembled full
sparse-matrix residual.
The candidate starts from the FFT-preconditioned conjugate-gradient result at relative tolerance
10\ensuremath{-}8, absolute tolerance 10\ensuremath{-}11 and maximum 512 iterations. It permits at most four additional
phases of 128 iterations, reevaluating the matrix-free residual between phases against a target 10\ensuremath{-}7.
These are bounded numerical-transfer controls, not a change to a production default.
4

Admission requires inclusive limits of 10\ensuremath{-}7 on the fresh matrix-free residual, 2 \ensuremath{\times} 10\ensuremath{-}5 on the
independently assembled full sparse residual, 10\ensuremath{-}10 on the compatible direct-reference residual, and
2\ensuremath{\times}10\ensuremath{-}4 on each gradient comparison. “Independent” here describes the original assembled-reference
path. Our present examination checks retained values and definitions; it does not recompute Ap \ensuremath{-} b
from recovered graph and solution arrays.
Shape; contrast Fresh rrel grel, compatible Base pass Candidate
8 \ensuremath{\times} 9 \ensuremath{\times} 10; 102 9.4084 \ensuremath{\times} 10\ensuremath{-}8 8.4343 \ensuremath{\times} 10\ensuremath{-}7 Yes Admit
8 \ensuremath{\times} 9 \ensuremath{\times} 10; 103 8.1041 \ensuremath{\times} 10\ensuremath{-}8 3.3308 \ensuremath{\times} 10\ensuremath{-}6 Yes Admit
8 \ensuremath{\times} 9 \ensuremath{\times} 10; 104 8.3258 \ensuremath{\times} 10\ensuremath{-}8 8.7205 \ensuremath{\times} 10\ensuremath{-}5 No Admit
10 \ensuremath{\times} 11 \ensuremath{\times} 12; 103 7.6658 \ensuremath{\times} 10\ensuremath{-}8 4.6783 \ensuremath{\times} 10\ensuremath{-}6 No Admit
Table 3: Four retained manufactured-graph candidate admissions. “Base pass” is the original per-
case aggregate outcome, not a relabeling of the candidate. Both failed baseline outcomes remain
part of the result.
All four candidates meet the captured numerical limits, including the two cases whose base compar-
isons fail. The aggregate base-panel pass therefore remains false. Neither the candidate admissions
nor the very small compatible direct-reference residuals establish physical calibration, application
accuracy or robustness beyond these four graphs. They establish a bounded algebraic acceptance
result under the captured definitions.

\section*{5 Couette startup: observable agreement preserves the physical}

failure

\section*{5.1 Case and observation definitions}

The annular case has inner radius ri = 3 m, outer radius ro = 7 m and axial length Lz = 19 .2 m.
The inner cylinder reaches Ωi = 0 .01 rad s \ensuremath{-}1 through a two-second smoothstep ramp; the outer
cylinder remains stationary. A 963 lattice uses \ensuremath{\Delta}x = 0 .2 m, \ensuremath{\Delta}t = 0 .04 s, \ensuremath{\tau} = 0 .8 and reference
density \ensuremath{\rho}0 = 1 kg m \ensuremath{-}3. There are 720 evolution steps and 721 saved states, including the initial
state, ending at 28.8 s. The physical viscosity used by the observer is
\ensuremath{\mu} = \ensuremath{\rho}0
(\ensuremath{\tau} \ensuremath{-} 1/2)\ensuremath{\Delta}x2
3\ensuremath{\Delta}t = 0.1 Pa s. (5)
These are stated numerical-to-physical scales, not a calibration against an apparatus.
The configuration uses 40 immersed-boundary iterations, an exterior moving-least-squares stencil
and post-stream reconstraint. Its physical marker interpolation is periodic axially, while the pop-
ulation core rolls periodically in all three Cartesian directions with the radial walls inside the box.
Native evolution is float32. Products and moment reductions are retained as float64; enabling a
64-bit runtime flag does not make the evolved population state double precision. The local diagonal
correction used in the periodic campaign is disabled here.
The steady annular reference profile is
u\ensuremath{\theta},s(r) = ar + b
r , a = Ωor2
o \ensuremath{-} Ωir2
i
r2o \ensuremath{-} r2
i
, b = r2
i r2
o(Ωi \ensuremath{-} Ωo)
r2o \ensuremath{-} r2
i
. (6)
5

The transient comparisons use a retained startup reference, not this steady curve at every time.
For angular velocity \ensuremath{\omega} = u\ensuremath{\theta}/r, that reference solves the axisymmetric incompressible diffusion
problem r3\ensuremath{\partial}t\ensuremath{\omega} = \ensuremath{\nu} \ensuremath{\partial}r(r3\ensuremath{\partial}r\ensuremath{\omega}), with \ensuremath{\nu} = \ensuremath{\mu}/\ensuremath{\rho}0. Fluid starts at rest; the inner-wall angular speed
follows the specified two-second cubic-smoothstep ramp and the outer wall remains stationary.
The axial interpretation is periodic or infinite-cylinder flow with no end plates. Body-fitted radial
control volumes are advanced with SciPy BDF. The retained profile reference uses 256 radial
intervals and the control-surface reference uses 512; both primary integrations use relative tolerance
10\ensuremath{-}9, angular-speed absolute tolerance 10\ensuremath{-}12 rad s\ensuremath{-}1, and maximum time step 0.125 s. These are
historical reference settings. The reference passed its own recorded admission checks but has not
been independently re-integrated in the present analysis.
Two observation sets must be distinguished. The primary surface observer uses z-averaged velocity,
density-weighted velocity, and second moments from all saved times. The velocity-profile check
instead averages the central two z planes at each of 1564 Cartesian sample locations satisfying
4 \ensuremath{\leq} r \ensuremath{\leq} 6 m, retaining angular duplicates. Seven explicit central-plane profiles survive at steps
0, 25, 50, 100, 200, 400, 720, together with all 721 per-time norm summaries. A full central-plane
trajectory cannot be reconstructed from z-averaged moments.

\section*{5.2 Independent reconstruction of the saved observer}

We independently implement the observer reductions without importing the original solver. Fourth-
order centered Cartesian derivatives and bilinear surface interpolation give the Newtonian shear
component on a cylindrical surface,
\ensuremath{\sigma}xx = 2\ensuremath{\mu} \ensuremath{\partial}xux, \ensuremath{\sigma}yy = 2\ensuremath{\mu} \ensuremath{\partial}yuy, \ensuremath{\sigma} xy = \ensuremath{\mu}(\ensuremath{\partial}yux + \ensuremath{\partial}xuy),
\ensuremath{\sigma}\ensuremath{\theta}r = cos \ensuremath{\theta} sin \ensuremath{\theta}(\ensuremath{\sigma}yy \ensuremath{-} \ensuremath{\sigma}xx) + (cos2 \ensuremath{\theta} \ensuremath{-} sin2 \ensuremath{\theta})\ensuremath{\sigma}xy,
Q(r, t) = r2Lz
Z 2\ensuremath{\pi}
0
\ensuremath{\sigma}\ensuremath{\theta}r(r, \ensuremath{\theta}, t) d\ensuremath{\theta}.
(7)
The primary reconstruction uses 256 angular points and four-point radial Gaussian quadrature for
volume observables. Pressure has zero axial torque on a concentric cylindrical surface by geometry.
This identity is not an independently measured pressure-accuracy result, and the reconstructed
stress is not an exact lattice momentum-flux tensor.
Density-weighted products are formed before averaging over z. For example, the angular momentum
of an annular volume is reduced from
LV = \ensuremath{\rho}0
Z
V
[x \ensuremath{\rho}reluy \ensuremath{-} y \ensuremath{\rho}relux] dV, (8)
with a separate constant-reference-density channel formed from u. The bar denotes the saved
axial mean and \ensuremath{\rho}rel the dimensionless lattice density. Replacing the average of a product with the
product of averages would change the observable. The same distinction applies to the convective
second moments.
All 721 saved times were reduced. The reconstructed outputs agree with the original 14 channels:
13 numeric channels and one recomputed finite-status predicate. One numeric channel is the zero
pressure-torque identity. They are therefore neither 14 independent experiments nor 14 independent
empirical measurements. The maximum viscous-torque difference is 1.33 \ensuremath{\times} 10\ensuremath{-}15 N m; a density-
weighted angular-momentum channel differs by at most 2.84 \ensuremath{\times} 10\ensuremath{-}14 kg m 2 s\ensuremath{-}1. These unit-specific
6

differences describe agreement of retained arithmetic, not physical error against the continuum
solution.
Doubling the angular torque sampling to 512 points gives a maximum difference of 5.96 \ensuremath{\times} 10\ensuremath{-}5
relative to the steady torque scale. The original report also contains doubled radial-volume and
reference-grid admission results. Our independent reconstruction does not repeat the complete
doubled radial-volume calculation or the original three-dimensional force-footprint audit. Effec-
tive loads inferred from storage and flux balances remain control-volume quantities, not directly
reconstructed wall Cauchy traction.

\section*{5.3 Physical acceptance remains failed}

At each control radius, the surface-stress error is normalized by the retained positive steady torque
scale Qs = 2.6600493 N m:
eQ(r) = max
t
|Q(r, t) \ensuremath{-} Qref (r, t)|
Qs
. (9)
This is a maximum over the startup history, not a comparison only at steady state. The three
control radii are 4.8, 5.0 and 5.2 m. Their errors are shown in Table 4. All exceed the captured 5\%
limit despite the close arithmetic reconstruction.
For the central-plane profile let ⃗ vj and ⃗ vref,j be the sampled azimuthal velocities. The two reported
profile endpoints are
eagg =
" P
j\ensuremath{\in}T \ensuremath{\parallel}⃗ vj \ensuremath{-} ⃗ vref,j\ensuremath{\parallel}2
2P
j\ensuremath{\in}T \ensuremath{\parallel}⃗ vref,j\ensuremath{\parallel}2
2
\#1/2
, e final = \ensuremath{\parallel}⃗ v720 \ensuremath{-} ⃗ vref,720\ensuremath{\parallel}2
\ensuremath{\parallel}⃗ vref,720\ensuremath{\parallel}2
, (10)
where T comprises the retained time observations with nonzero reference norm; the zero initial
reference does not contribute to the denominator. The aggregate is reconstructed from all 721
norm summaries, not from only the seven explicit profiles. Independently reducing those seven
profiles supports their pointwise comparisons and the final value, but does not authenticate absent
central-plane arrays at the other times.
Observable Error (\%) Limit (\%) Outcome
Surface torque, r = 4.8 m 7.06174 5 Fail
Surface torque, r = 5.0 m 6.17514 5 Fail
Surface torque, r = 5.2 m 5.52880 5 Fail
Aggregate central-plane profile 7.32722 5 Fail
Final central-plane profile 5.19725 5 Fail
Table 4: Retained Couette physical acceptance failures. Surface errors use the steady torque
scale over the complete saved history; profile errors use the transient-reference norms in Eq. ( 10).
Successful observer reconstruction does not alter these outcomes.
The original finite-state, density, mass, Mach-number, support and wall-slip checks pass. The
internal shell-balance diagnostic and observer-admission result also pass, while the surface-stress
and profile checks fail. Thus the saved trajectory is finite and internally checkable but does not
meet the specified physical-response tolerances.
Definitions, source bindings and corrected audit history are retained in the private reproducibility
record.
7


\section*{6 S-bend reference sensitivity: balance and convergence differ}

The separate curved-flow campaign uses OpenFOAM pimpleFoam at fixed flow 8.0424772 \ensuremath{\times} 10\ensuremath{-}7
m3 s\ensuremath{-}1. The retained coarse, medium and fine meshes contain 8471, 55919 and 128345 cells. Each
fixed-flow record ends after 512 steps at 20.48 s. Surface summaries retain inlet and outlet pressure
values, inlet and outlet fluxes, and net and absolute wall flux. The pressure-field dimensions
and density binding have not been recovered, so pressure differences are reported in their native
convention without conversion to Pa.
For two configurations, the pressure and balance comparisons are
\ensuremath{\Delta}p = pin \ensuremath{-} pout, \ensuremath{\delta} p = 100 \ensuremath{\Delta}pnew \ensuremath{-} \ensuremath{\Delta}pold
\ensuremath{\Delta}pold
, \ensuremath{\varepsilon} \ensuremath{\varphi} = |\ensuremath{\varphi}in + \ensuremath{\varphi}out + \ensuremath{\varphi}wall,net|
|\ensuremath{\varphi}in| . (11)
A common fixed pressure scale cancels from \ensuremath{\delta}p; this does not establish that scale’s physical units.
The flux signs follow the retained outward-normal convention.
Configuration Cells Native \ensuremath{\Delta}p \ensuremath{\varepsilon}\ensuremath{\varphi}
Coarse fixed flow 8,471 2.31759072 \ensuremath{\times} 10\ensuremath{-}4 3.1334 \ensuremath{\times} 10\ensuremath{-}11
Medium fixed flow 55,919 2.31357443 \ensuremath{\times} 10\ensuremath{-}4 1.1153 \ensuremath{\times} 10\ensuremath{-}10
Fine fixed flow 128,345 2.32634522 \ensuremath{\times} 10\ensuremath{-}4 1.7718 \ensuremath{\times} 10\ensuremath{-}10
Fine, four outer correctors 128,345 2.31199835 \ensuremath{\times} 10\ensuremath{-}4 1.3578 \ensuremath{\times} 10\ensuremath{-}10
Medium, sixteen outer correctors 55,919 2.30528250 \ensuremath{\times} 10\ensuremath{-}4 1.1352 \ensuremath{\times} 10\ensuremath{-}10
Table 5: Retained S-bend surface results. The fine-grid corrector variant is paired with the fine
fixed-flow case and increases outer correctors from two to four. The sixteen-corrector medium row
belongs to a separate eight-to-sixteen study; its eight-corrector baseline is not the listed medium
fixed-flow row. Corrector variants are iteration changes, not additional mesh levels.
The pressure sequence decreases from coarse to medium and then increases to fine. The medium-to-
fine relative change is +0.551994\%. On the fine mesh, increasing outer correctors from two to four
changes pressure difference by \ensuremath{-}0.616713\%, a comparable magnitude. The surface-flux imbalance
is at most 1.78 \ensuremath{\times} 10\ensuremath{-}10 across the retained variants, and the reported wall fluxes are zero. Thus
a highly balanced surface flux does not remove the pressure observable’s sensitivity to iteration
settings.
These records do not establish an asymptotic refinement regime. Nonmonotone pressure values,
nonuniform meshes and cell counts alone do not support a Richardson order or a grid-convergence
index. The corrector sensitivity cannot be subtracted as an independent mesh-error estimate with-
out a matched study. In the corrector variants, an unsupported residual-function warning accom-
panies a record that the residual monitor did not execute. It is not a measurement of a very small
residual.
An earlier common-sample postprocessing campaign refers to different solver runs. We therefore
do not join its velocity and near-wall comparisons to this fixed-flow mesh sequence as if they were
matched fields. Temporal qualification and an application-specific tolerance are also unresolved.
The S-bend evidence is an independent solver-family sensitivity example, not a continuum truth
against which the canonical JAX cases have been validated.
8


\section*{7 Implications and limits of the verification ladder}

The four campaigns answer different questions. Periodic decay checks a known smooth mode
under periodic conditions. Manufactured pressure admission checks a discrete algebraic problem
with a specified null-space and residual convention. Saved-state reconstruction checks whether an
observation algorithm consistently reduces its input. Couette startup checks a physical velocity and
stress response against the retained continuum reference. S-bend refinement examines sensitivity
to mesh and iteration choices. Passing one endpoint does not logically establish another endpoint
whose boundary conditions, observable or denominator differ.
The Couette reconstruction makes this distinction concrete. Very close agreement between two
reductions rules out some observation-arithmetic discrepancies for the retained moments; it does
not reduce the 5.53– 7.06\% surface-stress disagreement with the reference. Likewise, an internal shell
balance concerns consistency of storage and resolved fluxes, while a profile comparison concerns
agreement with a reference response. Neither requires the other to pass. The negative result is
useful precisely because successful internal checks coexist with failed physical tolerances.
The evidence does not isolate why the Couette response fails. The periodic and annular cases use
different boundaries and different correction flags. They also exercise different observables and
forcing conventions. Diagnosing a causal boundary mechanism would require matched numerical
interventions. Nothing here attributes the failure to a particular arithmetic contraction mode or
validates a separate rheology, mixer or stirred-tank application.
There are bounded reproducibility gaps. Native population trajectories and pressure graph/solution
vectors are absent, preventing independent evolution and full residual remeasurement. The tran-
sient annular reference has not been independently re-solved, and only seven complete central-plane
profiles are retained. Three-dimensional force footprints and the full doubled volume observer are
not independently reconstructed. S-bend pressure units, matched raw velocity fields and tempo-
ral convergence remain unqualified. These gaps restrict the interpretation; they do not erase the
physical failures that the available observations support.

\section*{8 Conclusion}

The retained campaigns support positive periodic decay and four manufactured-pressure candidate
admissions, alongside a failed annular physical response and an incompletely qualified curved-flow
reference. Independent reconstruction of all saved Couette observer outputs preserves the stress
and profile failures. In the S-bend records, small surface-flux imbalance coexists with nonmono-
tone mesh behavior and a material corrector sensitivity. The practical result is a set of distinct
acceptance decisions with explicit configurations and denominators. It supports further targeted
verification without treating successful internal arithmetic or algebraic stopping as completed phys-
ical qualification.
Data and reproducibility
Original source records, source-bound saved arrays, analysis definitions, independent reduction out-
puts and correction histories are privately retained in Supplementary record S1 with a provenance
index. They have not been publicly released, and external reviewer access and redistribution ar-
rangements are not established here. The separate unresolved-check report identifies the precise
9

additional inputs required for broader reproduction. No new fluid evolution was performed for this
retrospective analysis.

\section*{References}

[1] M. Geier, M. Schönherr, A. Pasquali and M. Krafczyk, “The cumulant lattice Boltzmann equa-
tion in three dimensions: Theory and validation,” Computers and Mathematics with Applica-
tions 70 (2015), 507–547. doi:10.1016/j.camwa.2015.05.001.
10

\end{document}
